\subsection{Attendus du projet, livrables et structuration du PMI}

Cette section a pour objectif de clarifier explicitement les attendus du Projet Master Innovation, les livrables associés, ainsi
que la manière dont le document est structuré afin de répondre de façon cohérente et progressive au sujet proposé. Elle
constitue un point de verrouillage méthodologique du rapport, en définissant précisément les objectifs poursuivis et le
périmètre de l’étude.

\subsubsection{Objectifs et périmètre de l’étude}

Le présent \acrshort{PMI} vise à évaluer la faisabilité technique et la pertinence expérimentale d’un système de séparation
d’étages de fusée reposant sur le principe de l’aérofreinage, dans le contexte d’un lanceur expérimental. L’objectif n’est pas
de proposer une solution industrialisable, mais de déterminer si une telle expérience peut être mise en œuvre dans des conditions
réalistes et conduire à l’obtention de données exploitables.\\

Dans ce cadre, le travail mené s’articule autour de plusieurs livrables principaux :
\begin{itemize}
    \item la conception et la réalisation de \textbf{prototypes fonctionnels} des systèmes de séparation et d’aérofreinage,
    intégrant leurs dispositifs de commande électroniques ainsi que la chaîne de mesure associée ;
    \item la production de \textbf{résultats de simulations numériques} permettant de caractériser les efforts
    aérodynamiques générés par les aérofreins et leur influence sur le comportement du lanceur ;
    \item la réalisation d’\textbf{analyses structurelles} visant à vérifier la tenue mécanique des éléments critiques et la
    cohérence globale de l’architecture retenue ;
    \item une \textbf{analyse globale des performances attendues}, des limites identifiées et de la pertinence du principe de
    séparation par aérofreinage dans un cadre expérimental.\\
\end{itemize}

\subsubsection{Structuration du travail et du document}

Le projet SP-02 mobilise des compétences variées en mécanique, aérodynamique, électronique et instrumentation, ce qui a conduit
à une organisation du travail fondée sur la spécialisation des membres selon ces domaines. La structure du document ne reflète
toutefois pas directement cette répartition interne, mais vise à proposer une lecture progressive et cohérente du projet. Les
différentes contributions y sont ainsi articulées autour d’un raisonnement global, depuis la conception des systèmes jusqu’à
l’analyse des résultats et aux réalisations.\\

La partie 2 est consacrée à la conception des systèmes de séparation et d’aérofreinage, ainsi qu’aux fonctions de séquencement
et de commande associées. Elle vise à présenter les choix architecturaux retenus et à définir les systèmes étudiés, tant du
point de vue mécanique qu’électronique.\\

La partie 3 présente la méthodologie employée pour étudier le fonctionnement des systèmes conçus. Elle regroupe les hypothèses
de modélisation, les approches numériques retenues pour les études aérodynamiques et structurelles, ainsi que les éléments
relatifs au processus expérimental envisagé. Cette partie a pour objectif de préciser le cadre d’analyse dans lequel les
résultats sont obtenus.\\

La partie 3 est consacrée à la méthodologie d’étude et de validation des systèmes conçus. Elle regroupe, au sein d’un même
ensemble cohérent, les hypothèses de modélisation retenues pour les études aérodynamiques et structurelles, ainsi que la
définition de la chaîne de mesure expérimentale envisagée pour la validation physique de l’expérience. Cette partie vise à
expliciter le cadre dans lequel les analyses numériques sont menées et à montrer comment elles s’articulent avec les moyens de
mesure prévus afin d’aboutir à des données exploitables.\\

Enfin, la partie 5 regroupe l’ensemble des réalisations concrètes issues du projet, incluant la fabrication des systèmes
mécaniques, le développement des systèmes électroniques et logiciels, ainsi que leur intégration. Cette partie permet de relier
les études menées aux éléments effectivement produits et d’apprécier la cohérence entre conception, analyses et mise en œuvre.\\

L’enjeu du PMI ne réside pas uniquement dans la réalisation d’études prises isolément, mais dans la capacité à mettre en regard
les différents résultats obtenus. Les choix de conception, les hypothèses aérodynamiques et les analyses structurelles sont
ainsi articulés afin d’éclairer progressivement le comportement du système étudié. La structuration du document vise à
accompagner cette démarche, en proposant une lecture cohérente permettant d’apprécier la faisabilité de l’expérience de
séparation par aérofreinage et la qualité des données susceptibles d’en être issues.

\newpage